% !TEX root = ../main.tex
\chapter{Apple Accelerate: BLAS, LAPACK, vDSP, vForce, BNNS}
\label{ch:accelerate}

The previous chapter described \MLX, the GPU-side substrate on which
\Lucid\ depends. This chapter describes its CPU-side counterpart:
\Accelerate, Apple's umbrella framework for high-performance
numerical computing on macOS and iOS \cite{apple_accelerate}.
\Accelerate\ has shipped with every release of macOS since 10.3
(Panther, 2003); it is a stable, documented, and unavoidable
dependency for any framework that wants to do serious CPU numerics
on Apple hardware. \Lucid's CPU backend
is built entirely on \Accelerate; no other CPU numerics library is
linked into the engine, and no \Lucid\ CPU code path hand-rolls a
SIMD kernel that \Accelerate\ already provides.

The choice is not purely pragmatic. \Accelerate\ is the only way to
access the AMX coprocessor (\chref{ch:apple_silicon}
Section~\ref{sec:apple_silicon:amx}) without resorting to
reverse-engineered instruction encodings, and the AMX path is the
difference between roughly $128$ GFLOP/s of NEON throughput per
P-core and roughly $2$ TFLOP/s of cluster-shared throughput on
matrix multiplication. To bypass \Accelerate\ would be to give up
that ten-fold gap. This chapter explains the sublayers of
\Accelerate, how \Lucid\ uses each, and what compatibility
considerations the framework has accumulated from the SDK churn of
the 2020--2024 period.

\section{Overview}
\label{sec:accelerate:overview}

\Accelerate\ is not a single library but an umbrella of related
libraries, each presented to client code as a top-level header but
backed by a shared dynamic library. The relevant sublayers, in the
order they were introduced, are:

\begin{itemize}
  \item \textbf{BLAS} (Basic Linear Algebra Subprograms): the
    Fortran-derived dense linear algebra interface, with Apple's
    implementation since macOS 10.3.
  \item \textbf{LAPACK} (Linear Algebra Package): the higher-level
    decomposition and solver suite, also shipped since 10.3.
  \item \textbf{vDSP} (Vector Digital Signal Processing): per-element
    vector operations (arithmetic, comparisons, packing/unpacking,
    FFT). Predates BLAS in the original Velocity Engine framework
    on PowerPC.
  \item \textbf{vForce}: transcendental and special functions on
    vectors. A separate header for historical reasons (vForce was
    originally a separate library; the symbols now live in
    \code{libvecLib.dylib} like the rest of \Accelerate).
  \item \textbf{vImage}: image-processing primitives. Not used by
    \Lucid; mentioned for completeness.
  \item \textbf{BNNS} (Basic Neural Network Subroutines): a
    higher-level convolution and pooling library introduced in
    macOS 10.13. The first \Accelerate\ sublayer with an explicit
    ML focus.
  \item \textbf{SparseSolvers}: sparse linear algebra. Not used by
    \Lucid; sparse tensors are out of scope for \Lucid\ 3.x.
\end{itemize}

The four sublayers \Lucid\ actually uses are BLAS, LAPACK, vDSP,
vForce, and BNNS. We treat each in turn.
\tabref{tab:accelerate:sublayers} gives a one-page overview before
we dive into each.

\begin{table}[t]
\centering
\small
\begin{tabular}{lllll}
\toprule
Sublayer & Introduced & Path to silicon & Lucid usage & File \\
\midrule
BLAS    & 10.3 (2003)  & AMX coprocessor   & GEMM, axpy, dot     & \texttt{AccelerateMath.cpp} \\
LAPACK  & 10.3 (2003)  & BLAS + LAPACK     & cholesky, qr, svd   & \texttt{AccelerateMath.cpp} \\
vDSP    & 10.0 (PowerPC) & NEON SIMD       & elementwise, reductions & \texttt{AccelerateMath.cpp} \\
vForce  & 10.4 (2005)  & SIMD + lookup     & exp, log, sin, erf  & \texttt{AccelerateMath.cpp} \\
BNNS    & 10.13 (2017) & AMX / SME         & conv2d, pooling     & \texttt{AccelerateConv.cpp} \\
vImage  & 10.3 (2003)  & SIMD + custom     & \emph{not used}     & --- \\
SparseSolvers & 10.13  & sparse BLAS       & \emph{not used}     & --- \\
\bottomrule
\end{tabular}
\caption[Accelerate sublayers and \Lucid\ usage.]{The seven
sublayers of Apple's \Accelerate\ framework, by introduction date,
the hardware path each ultimately reaches, the operations \Lucid\
delegates to each, and the \Lucid\ source file that wraps it. The
two unused sublayers (vImage, SparseSolvers) are listed for
completeness.}
\label{tab:accelerate:sublayers}
\end{table}

\subsection{Linkage and headers}
\label{ssec:accelerate:linkage}

\Accelerate\ is exposed to C and C++ code through the umbrella header
\code{<Accelerate/Accelerate.h>}, which transitively includes
\code{<vecLib/vecLib.h>}, \code{<vecLib/cblas.h>},
\code{<vecLib/clapack.h>}, and so on. The umbrella links the
dynamic library \code{Accelerate.framework}, which is part of the
macOS base system and present in every install.

\Lucid's CMake build adds \code{Accelerate.framework} to the link
line for the engine target. There is no separate find-package step
because the framework is in the standard SDK location and
\code{find\_library(ACCELERATE\_FRAMEWORK Accelerate)} resolves it
unconditionally on Apple silicon.

\subsection{Two SDK transitions}
\label{ssec:accelerate:sdk_transitions}

Two SDK transitions in the 2022--2024 period have affected the
\Lucid\ build:

\begin{enumerate}
  \item \textbf{\code{ACCELERATE\_NEW\_LAPACK}.} In macOS 13.3,
    Apple introduced a new LAPACK ABI that uses 64-bit integers for
    matrix dimensions, replacing the legacy 32-bit ABI inherited
    from netlib LAPACK. To opt into the new ABI, client code must
    define \code{ACCELERATE\_NEW\_LAPACK=1} and
    \code{ACCELERATE\_LAPACK\_ILP64=1} before including the headers.
    \Lucid\ adopted the new ABI in version 3.0.2 as part of the
    numpy-independence pass (vault:
    \code{retro-3-0-2-numpy-free-standalone}), because the new ABI
    is the only one that interoperates correctly with large
    matrices.
  \item \textbf{BNNS deprecation warnings.} Starting with macOS 14,
    several BNNS functions are marked \code{API\_DEPRECATED} in the
    public headers, on the grounds that CoreML provides equivalent
    or higher-level functionality. The functions still work and are
    still the fastest CPU convolution path; \Lucid\ silences the
    warning by passing \code{-Wno-deprecated-declarations} only to
    the file that includes the BNNS header, not framework-wide. The
    discipline is documented in \code{arch-device-bifurcation-complete}.
\end{enumerate}

\section{BLAS and LAPACK}
\label{sec:accelerate:blas}

BLAS and LAPACK are the workhorses of \Lucid's CPU linear algebra.
Together they cover matrix multiplication, factorisations
(Cholesky, QR, SVD, eigendecomposition), and linear-system solvers.
\Lucid\ uses them through the CBLAS and CLAPACK C-language
interfaces, which avoid the Fortran calling convention but preserve
the underlying column-major storage convention.

\subsection{The CBLAS interface}
\label{ssec:accelerate:cblas}

CBLAS exposes the BLAS routines as C functions with a layout
argument that selects between \code{CblasRowMajor} and
\code{CblasColMajor}. \Lucid's CPU backend uses \code{CblasRowMajor}
throughout, matching \Lucid's tensor storage convention. The
function pattern is:

\begin{lstlisting}[style=lucidcpp,language=C++]
#include <Accelerate/Accelerate.h>

// C := alpha * A * B + beta * C
cblas_sgemm(
    CblasRowMajor,
    CblasNoTrans, CblasNoTrans,
    M, N, K,            // dims
    alpha,
    A_ptr, K,           // lda = stride of leading dimension
    B_ptr, N,           // ldb
    beta,
    C_ptr, N);          // ldc
\end{lstlisting}

The \code{cblas\_sgemm} call, with operands large enough to amortise
dispatch overhead, executes on the AMX coprocessor and achieves
the throughput characterised in
\chref{ch:apple_silicon}~Section~\ref{ssec:apple_silicon:amx_perf}.
For operands below roughly $64 \times 64$, the AMX path is not
selected (the dispatch overhead exceeds the win), and the call
runs through NEON; for very small operands, scalar code in
\Lucid's composite layer is sometimes faster still.

The double-precision counterpart is \code{cblas\_dgemm}; \Lucid\
exposes it but does not exercise it on the GPU path, because
\Metal\ has no FP64 support and a mixed-device FP64 GEMM is rare
in practice.

\subsection{LAPACK routines used by Lucid}
\label{ssec:accelerate:lapack_used}

\Lucid's \code{linalg} subpackage (\chref{ch:linalg}) wraps the
following LAPACK routines, in their CLAPACK form with the new
ILP64 ABI:

\begin{itemize}
  \item \code{cholesky} \(\to\) \code{spotrf} / \code{dpotrf}
    (Cholesky factorisation, lower-triangular).
  \item \code{qr} \(\to\) \code{sgeqrf} + \code{sorgqr}
    (QR factorisation followed by Q reconstruction).
  \item \code{svd} \(\to\) \code{sgesdd} (divide-and-conquer SVD;
    asymptotically faster than \code{sgesvd}).
  \item \code{eigh} \(\to\) \code{ssyevd} (divide-and-conquer
    symmetric eigendecomposition).
  \item \code{solve} \(\to\) \code{sgesv} (general LU-based
    solve).
  \item \code{lstsq} \(\to\) \code{sgelsd} (divide-and-conquer
    least-squares).
  \item \code{inv} \(\to\) \code{sgetrf} + \code{sgetri} (LU
    followed by explicit inverse construction).
  \item \code{det}, \code{logdet} \(\to\) \code{sgetrf} followed
    by Python-side product of diagonal entries.
\end{itemize}

The choice of divide-and-conquer variants (\code{sgesdd},
\code{ssyevd}, \code{sgelsd}) over the simpler QR-based variants is
deliberate: for matrices larger than a few hundred entries on each
side, divide-and-conquer is significantly faster and Apple's
implementation in \Accelerate's LAPACK is competitive with the best
open-source LAPACK builds.

\subsection{The linalg carve-out, revisited}
\label{ssec:accelerate:linalg_carveout}

\chref{ch:mlx_runtime}~Section~\ref{ssec:mlx_runtime:linalg_carveout}
introduced the linalg carve-out: even when the user has selected
\code{device=`metal'}, linear-algebra primitives dispatch on the CPU
stream, because \Metal\ has no LAPACK-equivalent kernels.

The mechanism is now visible from this chapter's perspective.
\MLX's CPU stream uses, internally, the same \Accelerate\ BLAS and
LAPACK that \Lucid's CPU backend uses directly. When \Lucid\
dispatches a \code{cholesky} on a \code{metal} tensor, the path is:

\begin{enumerate}
  \item \Lucid\ moves the tensor through \code{SharedStorage} to a
    CPU view (this is the bridge of \chref{ch:unified_memory}).
  \item \Lucid\ calls \code{cblas\_sgemm}-style \Accelerate\ LAPACK
    on the CPU view. Under the hood, \Accelerate\ uses the AMX
    coprocessor for BLAS-3 portions and CPU code for the rest.
  \item The result, residing in unified memory, is wrapped as an
    \MLX\ array and returned as a metal-device tensor.
\end{enumerate}

The bridge is free in the contiguous case
(\chref{ch:unified_memory}~Section~\ref{ssec:unified_memory:bridge_free});
the only real cost is the LAPACK computation itself, which is
several orders of magnitude larger than the bridge overhead for any
matrix worth decomposing.

\section{vDSP}
\label{sec:accelerate:vdsp}

vDSP is the lowest layer of \Accelerate\ that \Lucid\ uses. It
provides per-element vector operations on contiguous arrays: scalar
multiplication and addition, dot products, element-wise arithmetic,
sums, means, and a small set of packing and conversion utilities.
The primitives are mostly C functions with a strict naming
convention: \code{vDSP\_<op>} for single-precision,
\code{vDSP\_<op>D} for double-precision.

\subsection{Primitives used by Lucid}
\label{ssec:accelerate:vdsp_used}

The CPU backend's elementwise kernels are thin wrappers over vDSP
primitives. The most-used:

\begin{itemize}
  \item \code{vDSP\_vmul} / \code{vDSP\_vadd} / \code{vDSP\_vsub} /
    \code{vDSP\_vdiv}: elementwise binary arithmetic on
    single-precision vectors.
  \item \code{vDSP\_vsadd} / \code{vDSP\_vsmul} / \code{vDSP\_vsdiv}:
    scalar-broadcast variants.
  \item \code{vDSP\_vneg}, \code{vDSP\_vabs}, \code{vDSP\_vsq}:
    unary arithmetic.
  \item \code{vDSP\_dotpr}, \code{vDSP\_sve} (sum), \code{vDSP\_meanv}
    (mean), \code{vDSP\_maxv}, \code{vDSP\_minv}: reductions.
  \item \code{vDSP\_vthr} (threshold), \code{vDSP\_vclip} (clamp).
  \item \code{vDSP\_vfix32} (float-to-int conversion),
    \code{vDSP\_vflt32} (int-to-float).
\end{itemize}

The full list \Lucid\ exposes is in
\code{lucid/\_C/backend/cpu/AccelerateMath.cpp}, where each is
wrapped behind a function with \refframework-compatible
naming.

\subsection{Stride support}
\label{ssec:accelerate:vdsp_strides}

vDSP primitives accept explicit stride arguments (the \code{N} and
\code{stride} pair after each pointer), which means that they can
operate on non-contiguous slices of a larger buffer without first
copying to contiguity. \Lucid's CPU kernels exploit this for views
that are unit-stride along the innermost dimension and arbitrarily
strided along the outer dimensions; the dispatcher emits a vDSP
call per outer-axis iteration.

For views that are non-unit-stride along the innermost dimension
(produced by certain transposes), \Lucid's CPU kernels first
materialise a contiguous copy. The trade-off is amortised by the
fact that such views are rare in practice and the contiguous copy
cost is small relative to the subsequent kernel.

\subsection{Why not OpenBLAS or MKL}
\label{ssec:accelerate:vdsp_why_not_other}

The natural alternative to vDSP would have been OpenBLAS, BLIS, or
the Intel MKL. All three are excellent libraries, but none of them
target AMX on Apple silicon. OpenBLAS on M-series uses NEON; BLIS
uses NEON; MKL does not run on ARM at all. Bypassing vDSP would
therefore have meant giving up the AMX path, which is the dominant
single performance source for CPU-side BLAS-3 on Apple silicon.

We have also benchmarked the OpenBLAS NEON path against \Accelerate\
vDSP plus AMX on representative GEMM sizes and \Accelerate\ is
faster by a factor of $3$--$5\times$ on $256 \times 256$ FP32 and
larger. The conclusion is conclusive.

\section{vForce}
\label{sec:accelerate:vforce}

vForce provides vectorised transcendental and special functions.
The interface convention matches vDSP but the functions are
mathematically richer:

\begin{itemize}
  \item \code{vvexpf}, \code{vvlogf}, \code{vvsinf}, \code{vvcosf},
    \code{vvtanf} and their double-precision \code{vv*} counterparts.
  \item \code{vvtanhf}, \code{vvsinhf}, \code{vvcoshf}.
  \item \code{vvsqrtf}, \code{vvrsqrtf}.
  \item \code{vverff}, \code{vverfinvf}.
\end{itemize}

\Lucid's CPU backend uses vForce wherever a primitive maps onto one
of these; the wrappers are in
\code{lucid/\_C/backend/cpu/AccelerateMath.cpp} alongside the vDSP
wrappers. The win over a scalar transcendental loop is roughly
$4$--$6\times$ on FP32 \code{exp} and \code{log}, falling to
$2$--$3\times$ on the more complex \code{erf} family.

A practical caveat: vForce's \code{vverfinvf} produces results that
differ in the last few ULPs from the scalar \code{erfinvf} that
ships in \code{libm.dylib}. \Lucid's gradcheck tolerance for
\code{erfinv}-using ops is relaxed accordingly
(\chref{ch:gradcheck}).

\section{BNNS}
\label{sec:accelerate:bnns}

BNNS is the highest level of \Accelerate\ that \Lucid\ uses, and the
most controversial. It is a convolution and neural-network primitive
library introduced in macOS 10.13 (2017), at a time when ML inference
on CPU was a more common need than it is now. Its API is shaped
around the standard NN-inference vocabulary: tensors with
explicit layout descriptors, convolution and pooling descriptors,
fully-connected layers, activation functions, and a small set of
loss-derivative primitives.

\subsection{The convolution fast path}
\label{ssec:accelerate:bnns_conv}

The single most consequential BNNS function for \Lucid\ is
\code{BNNSDirectApplyConvolutionLayer} (and its descriptor
constructor \code{BNNSDirectFilterCreateConvolutionLayer}). On
M-series hardware, this function dispatches the convolution kernel
through the AMX coprocessor, achieving throughput unattainable
through pure NEON code. For a $3 \times 3$ stride-1 padded
convolution on a $224 \times 224$ input with 64 output channels and
batch size 1, BNNS is roughly $8\times$ faster than a hand-written
NEON convolution loop and $3\times$ faster than calling
\code{cblas\_sgemm} on an im2col-flattened version.

\Lucid's CPU convolution kernels (\code{lucid/\_C/backend/cpu/AccelerateConv.cpp})
delegate to BNNS for the forward pass of \code{conv1d}, \code{conv2d},
and \code{conv3d}, with descriptor caching so that repeated
applications of the same layer to inputs of identical shape do not
re-pay the descriptor-construction cost.

\subsection{The deprecation warning}
\label{ssec:accelerate:bnns_deprecation}

Starting with macOS 14, several BNNS functions are marked
\code{API\_DEPRECATED} in the public headers. Apple's stated reason
is that CoreML provides equivalent or superior functionality at a
higher abstraction level. The deprecation does not affect the
runtime: the functions still work, still perform identically, and
are still part of the system framework. Only the compile-time
warning is generated.

\Lucid\ silences the warning narrowly, in the single file that
includes the BNNS header:

\begin{lstlisting}[style=lucidcpp,language=C++]
// In AccelerateConv.cpp
#pragma clang diagnostic push
#pragma clang diagnostic ignored "-Wdeprecated-declarations"
#include <Accelerate/Accelerate.h>
#pragma clang diagnostic pop
\end{lstlisting}

The framework-wide \code{-Wno-deprecated-declarations} flag is
explicitly \emph{not} used, because suppressing deprecation
warnings globally would mask genuine bug-class issues elsewhere.
The narrow scope is documented in \code{arch-device-bifurcation-complete}.

\subsection{Why Lucid still uses BNNS}
\label{ssec:accelerate:bnns_why}

Three reasons keep BNNS in \Lucid's CPU backend:

\begin{enumerate}
  \item It is the only public Apple-supported path to CPU-side
    matrix-coprocessor convolution. The alternative is reverse-
    engineered AMX assembly, which is unsupportable.
  \item The deprecation has not been followed by a removal. Apple
    has marked the functions deprecated since macOS 14 without
    removing them in macOS 15 or 26; we estimate the removal is
    several years away, at which point SME-based replacements will
    have matured.
  \item CoreML, the recommended replacement, requires offline graph
    compilation. \Lucid's CPU backend dispatches dynamically and
    cannot accept the offline-compilation requirement without
    losing the dynamic-shape support that the rest of the framework
    is built on.
\end{enumerate}

The plan is to retain BNNS until either (a) Apple removes the
functions and forces a migration, or (b) SME-based direct emission
from \code{lucid.compile} (\chref{ch:compile_design}) reaches parity
with BNNS throughput, whichever happens first.

\section{AccelerateBackend in Lucid}
\label{sec:accelerate:lucid_use}

The C++ class that mediates between \Lucid's dispatcher and
\Accelerate\ is \code{AccelerateBackend}
(\code{lucid/\_C/backend/cpu/AccelerateBackend.\{h,cpp\}}). It
implements the same \code{IBackend} interface as \code{MlxBackend}
(\chref{ch:mlx_runtime}~Section~\ref{ssec:mlx_runtime:mlx_backend}),
which is what allows the dispatcher to route operations to either
backend uniformly.

\subsection{File layout}
\label{ssec:accelerate:file_layout}

The backend is split across several files to manage the size of the
implementation:

\begin{itemize}
  \item \code{AccelerateBackend.\{h,cpp\}}: the \code{IBackend}
    implementation; dispatches per-op to the helpers below.
  \item \code{AccelerateMath.cpp}: vDSP and vForce wrappers for
    elementwise and reduction operations.
  \item \code{AccelerateConv.cpp}: BNNS wrappers for convolution
    forward; backward is composed from BLAS GEMM on im2col-flattened
    gradient tensors.
  \item \code{AccelerateBatchNorm.cpp}, \code{AccelerateLayerNorm.cpp},
    \code{AccelerateInstanceNorm.cpp}, \code{AccelerateGroupNorm.cpp}:
    normalisation kernels that fold mean / variance computation
    with the normalising affine in a single pass.
  \item \code{AccelerateRandom.cpp}: vDSP's random number generator
    wrapped for \Lucid's \code{Generator} interface.
  \item \code{Helpers.cpp}: tensor-to-vDSP shape and stride
    conversions, descriptor caching, layout adapters.
\end{itemize}

The split is the result of incremental growth across multiple
phases; the natural unit of factoring is the underlying
\Accelerate\ sublayer, not the abstract op kind, because each
sublayer has its own conventions and quirks.

\subsection{Per-op routing inside the backend}
\label{ssec:accelerate:per_op}

Inside \code{AccelerateBackend}, the per-op routing is a switch on
the abstract op kind:

\begin{lstlisting}[style=lucidcpp,language=C++]
TensorImpl AccelerateBackend::binary(
    OpKind op, const TensorImpl& a, const TensorImpl& b) {
  switch (op) {
    case OpKind::Add: return accel::vmul_vadd_path(a, b, /*op=*/Add);
    case OpKind::Mul: return accel::vmul_vadd_path(a, b, /*op=*/Mul);
    case OpKind::Pow: return accel::pow_path(a, b);
    // ... and so on
  }
}
\end{lstlisting}

The helpers in the \code{accel} namespace are the actual vDSP /
vForce wrappers; the backend class itself contains only routing.
This structure mirrors \code{MlxBackend}'s and is what allows the
two backends to be substitutable through the same dispatcher.

\section{Threading and Concurrency}
\label{sec:accelerate:threading}

\Accelerate\ is internally multithreaded for large operations. The
threading is managed by Grand Central Dispatch (GCD), the macOS
system-wide task scheduler. Client code does not configure thread
counts directly; instead, GCD allocates worker threads from its
shared pool based on system load and the operation's estimated
cost.

\subsection{The implicit thread pool}
\label{ssec:accelerate:threadpool}

This implicit pooling has two consequences for \Lucid:

\begin{enumerate}
  \item There is no \code{omp\_set\_num\_threads}-equivalent that
    the user can set. The number of threads \Accelerate\ uses for a
    given operation is determined at runtime by GCD.
  \item Operations from \Accelerate\ can compete with \Lucid's own
    Python-level threads (e.g.\ data-loader workers) for CPU
    cycles. The shared GCD pool prevents pathological
    over-subscription, but there is no explicit isolation between
    domains.
\end{enumerate}

In practice the scheduling is sensible: training workloads that
dispatch from a single Python thread to \Accelerate\ get
near-optimal P-core utilisation, and data-loading workers correctly
yield to the training thread under load. \Lucid\ does not set
explicit affinity hints
(\chref{ch:apple_silicon}~Section~\ref{ssec:apple_silicon:scheduling})
because the default behaviour is correct.

\subsection{Lock-free reentrancy}
\label{ssec:accelerate:reentrancy}

\Accelerate's functions are reentrant: it is safe for two threads to
call \code{cblas\_sgemm} simultaneously on disjoint inputs.
\Lucid's autograd engine relies on this to dispatch independent
operations from the optimiser step in parallel with the next
forward pass's preparation (\chref{ch:optimizers}).

\section{Quirks and Pitfalls}
\label{sec:accelerate:quirks}

We close with the three classes of issues that have surfaced during
\Lucid's development.

\subsection{LAPACK error returns}
\label{ssec:accelerate:lapack_errors}

LAPACK routines do not raise exceptions. They return an integer
status code through the \code{info} out-parameter; a negative value
indicates an invalid argument, a positive value indicates a numerical
failure (singular matrix in \code{getrf}, non-positive-definite in
\code{potrf}, non-convergence in \code{gesdd}).

\Lucid's wrappers check \code{info} after every call and translate
positive values into \code{LucidError(``cholesky failed: matrix not
positive definite at row $k$'')}-style messages. Negative \code{info}
values indicate a bug in the wrapper and are translated into
internal assertions, since they should be unreachable.

\subsection{Row-major vs.\ column-major}
\label{ssec:accelerate:row_col}

The native LAPACK convention is column-major; CLAPACK's
\code{LAPACKE\_*} wrappers accept a layout argument to allow
row-major input. \Lucid\ uses the CLAPACK \code{LAPACKE\_*} variants
throughout, with \code{LAPACK\_ROW\_MAJOR}. The conversion overhead
is internal to LAPACKE and is amortised over the decomposition cost.

A historical bug (3.0.3) involved calling the older non-LAPACKE
variant of \code{spotrf} and forgetting to transpose the input; the
result was a successful-looking Cholesky of the matrix's transpose.
The fix was to standardise on the LAPACKE wrappers throughout, and
the standardisation is now enforced by a code-review checklist
referenced in the \code{retro-3-0-3-perf-pass} note.

\subsection{SDK version drift}
\label{ssec:accelerate:sdk_drift}

\Accelerate's API surface is largely stable, but individual
functions have been added, deprecated, or had their behaviour
clarified across SDK versions. The pattern that has caused the
most churn for \Lucid\ is the LAPACK ILP64 ABI transition discussed
in \secref{ssec:accelerate:sdk_transitions}. The discipline now is:

\begin{enumerate}
  \item Pin a minimum macOS deployment target (\Lucid\ uses
    \code{MACOSX\_DEPLOYMENT\_TARGET=26.0}, matching \MLX's wheel).
  \item Use only documented APIs that have been stable across the
    minimum-target SDK and the current SDK.
  \item Treat any deprecation warning that surfaces during build as
    a TODO to either migrate to the suggested replacement or, if
    the replacement does not exist (BNNS being the canonical
    example), silence narrowly at the include site.
\end{enumerate}

The minimum-target choice is shared with \MLX\ because the two
dependencies' wheel tags must agree; \chref{ch:mlx_runtime}
~Section~\ref{ssec:mlx_runtime:macos_constraint} documents the
constraint from \MLX's side.

\section{Summary and Forward References}
\label{sec:accelerate:summary}

\Accelerate\ is the CPU-side substrate on which \Lucid\ stands. Its
sublayers --- BLAS, LAPACK, vDSP, vForce, BNNS --- jointly provide
matrix multiplication, decompositions, elementwise and
transcendental kernels, and convolution. The framework is the only
public Apple-supported path to the AMX coprocessor and is therefore
the only path that achieves competitive CPU-side ML throughput on
Apple silicon. \Lucid's \code{AccelerateBackend} wraps the
relevant subset behind an \code{IBackend} interface that the
dispatcher uses uniformly with \code{MlxBackend}.

The remaining hardware-foundation chapter, \chref{ch:metal_mpsgraph},
treats \Metal\ and \MPSGraph\ directly, since \Lucid's
graph-capture compiler emits onto \MPSGraph\ rather than relying on
\MLX's compile path. After that, the book moves into
\partref{part:system}, which formalises the layered architecture
the hardware chapters have implicitly sketched and explains how
\Lucid's tensor model, allocator, dispatch, and bridge boundaries
fit together at the engine level.

\begin{lucidtip}
For the user, the visible consequence of this chapter is that
\code{device=`cpu'} in \Lucid\ corresponds to AMX-accelerated BLAS
and BNNS through \Accelerate, not to a generic CPU fallback. This
makes CPU-only training feasible for small models and certain
non-convolutional workloads where the GPU's launch overhead would
dominate. For training-time benchmarks see \chref{ch:perf_methodology}.
\end{lucidtip}
